\section{Conclusion}
\label{sec:conclusion}

This paper treated early warning as a stateful alarm policy whose alert load is a declared design constraint rather than an emergent threshold artefact. We derived a patient-blocked finite-grid upper-confidence procedure for the population mean of patient-level alarms per day, separated it from a deliberately invalid timestep effective-sample-size shortcut, and implemented the complete path from causal scoring through policy replay to auditable operating points. On the public 2019 A/B cohort, the procedure produced a monotone internal budget-sensitivity frontier, with A sensitivity rising from 0.185 at 0.25 alarms per patient-day to 0.774 at 2 per day over three seeds. The same A-calibrated threshold did not transport to site B: at the 1-per-day budget, B sensitivity was 0.315 versus 0.585 at A, even though the realised B alarm rate remained below the A certificate. An oracle target-site recalibration raised B sensitivity to 0.497 but also raised the realised rate to 0.616, and a nested label-budget analysis showed that roughly 10\% of the eligible target subjects (about 1,482 patients) recovered the full-pool threshold at 1 per day while 5\% remained unstable across seeds. LightGBM raised B sensitivity at 1 per day to 0.540 and the GRU reached 0.375, but neither removed the cross-site shift; the timestep comparator was less conservative at tight budgets and increased patient-level budget violations. These results support budgeted, patient-blocked calibration as an engineering instrument, but they do not establish clinical effectiveness, alert-fatigue reduction or deployment readiness; label-free target adaptation, subgroup and missingness analyses, prospective outcome evaluation and a finer utility-aware grid remain necessary before any clinical claim.
